\documentclass[conference]{IEEEtran}

\usepackage{amsmath}
\usepackage{graphicx}
\usepackage{booktabs}
\usepackage{cite}

\title{Adaptive Multi-Modal SLAM with Uncertainty-Aware Sensor Fusion for Robust Autonomous Navigation}

\author{Panagiota Grosdouli}

\begin{document}

\maketitle

\begin{abstract}
Visual-inertial SLAM systems are central to autonomous robot navigation, yet their reliability can degrade under low illumination, motion blur, texture scarcity, and rapid motion. This paper investigates an uncertainty-aware adaptive fusion framework that estimates modality reliability online and adjusts the contribution of visual and inertial measurements during state estimation. The long-term goal is a self-healing SLAM system that can predict localization failure and activate recovery policies before catastrophic tracking loss.
\end{abstract}

\section{Introduction}
Autonomous robots operating in real-world environments must localize reliably despite perceptual degradation. Classical SLAM pipelines often assume fixed sensor fusion strategies, although the reliability of visual and inertial measurements changes over time.

\section{Related Work}
This section will review visual SLAM, visual-inertial odometry, uncertainty estimation, event-based sensing, and failure recovery in robotic perception.

\section{Method}
The proposed framework consists of a SLAM backend, uncertainty estimator, adaptive fusion policy, failure predictor, and recovery policy.

\section{Experiments}
Experiments will evaluate the system on EuRoC MAV and TUM-VI under controlled degradation.

\section{Results}
This section will report trajectory accuracy, tracking failures, recovery latency, and uncertainty calibration.

\section{Conclusion}
The project aims to demonstrate that explicit uncertainty reasoning can improve SLAM robustness and provide a path toward self-healing robot perception.

\bibliographystyle{IEEEtran}
\bibliography{references}

\end{document}
